\documentclass[12pt]{article}
\usepackage{amsmath, amssymb, amsthm}
\usepackage{geometry}
\geometry{margin=1in}

\theoremstyle{remark}
\newtheorem*{remark}{Remark}

\newcommand{\E}{\mathbb{E}}
\newcommand{\F}{\mathcal{F}}

\title{Piazza Response: Week 8, Exercise 4\\
       \large (for Jeremy Weissman)}
\author{Neil Hwang}
\date{}

\begin{document}
\maketitle

\section*{The Problem}

Let $X_1, X_2, \ldots$ be i.i.d.\ with $\E[X_1] = 0$ and $\E|X_1| < \infty$.
Set $S_n = X_1 + \cdots + X_n$ (with $S_0 = 0$) and
$\F_n = \sigma(X_1, \ldots, X_n)$.
Using Exercise~3 (``taking out what is known''), show that
\[
  \E[S_n \mid \F_{n-1}] = S_{n-1} \quad \text{a.s.\ for all } n \ge 1,
\]
and conclude that $(S_n, \F_n)_{n \ge 0}$ is a martingale.

\section*{Hints and Key Ideas}

There are exactly two things happening inside $\E[S_n \mid \F_{n-1}]$,
each for a different reason.  The trick is to separate them cleanly.

\paragraph{Step 1: Split the sum.}
Write $S_n = S_{n-1} + X_n$.  By \emph{linearity} of conditional expectation:
\[
  \E[S_n \mid \F_{n-1}]
  = \E[S_{n-1} \mid \F_{n-1}] + \E[X_n \mid \F_{n-1}].
\]
Now handle each term separately.

\paragraph{Step 2: First term --- ``taking out what is known.''}
The sum $S_{n-1} = X_1 + \cdots + X_{n-1}$ is determined entirely by
$X_1, \ldots, X_{n-1}$, so it is $\F_{n-1}$-measurable.  Apply Exercise~3
with $f = S_{n-1}$ and $g = 1$ (note $S_{n-1} \cdot 1 = S_{n-1} \in L^1$
since $\E|X_1| < \infty$, so the hypotheses are satisfied):
\[
  \E[S_{n-1} \mid \F_{n-1}]
  = \E[S_{n-1} \cdot 1 \mid \F_{n-1}]
  = S_{n-1} \cdot \E[1 \mid \F_{n-1}]
  = S_{n-1} \cdot 1
  = S_{n-1}.
\]
Here we used $\E[1 \mid \F_{n-1}] = 1$, since the conditional expectation of
any constant is that same constant (a constant is $\F_{n-1}$-measurable and
integrates to itself over every $A \in \F_{n-1}$).

\begin{remark}
  Alternatively, one can cite directly that if $f$ is $\F$-measurable then
  $\E[f \mid \F] = f$ a.s.\ (one of the standard properties of conditional
  expectation).  Applying Exercise~3 with $g = 1$ is one way to derive this
  from first principles, which is what the problem is asking for.
\end{remark}

\paragraph{Step 3: Second term --- independence.}
The increment $X_n$ is \emph{independent} of $\F_{n-1} = \sigma(X_1,\ldots,X_{n-1})$:
since the $X_i$ are i.i.d., $X_n$ is independent of any function of
$X_1, \ldots, X_{n-1}$.  By the independence property of conditional expectation
(if $W \perp \F$ then $\E[W \mid \F] = \E[W]$):
\[
  \E[X_n \mid \F_{n-1}] = \E[X_n] = \E[X_1] = 0.
\]

\paragraph{Combining.}
\[
  \E[S_n \mid \F_{n-1}]
  = \underbrace{S_{n-1}}_{\text{Step 2}} + \underbrace{0}_{\text{Step 3}}
  = S_{n-1} \quad \text{a.s.}
\]

\section*{Verifying the Three Martingale Conditions}

Having established the martingale property, you still need to check all three
conditions in the definition of a martingale
$(S_n, \F_n)_{n \ge 0}$:

\begin{enumerate}
  \item \textit{Adapted}: $S_n = X_1 + \cdots + X_n$ is a Borel function of
        $(X_1, \ldots, X_n)$, hence $\F_n$-measurable. $\checkmark$

  \item \textit{Integrable}: By the triangle inequality and the i.i.d.\ assumption,
        \[
          \E|S_n| \le \sum_{k=1}^n \E|X_k| = n\,\E|X_1| < \infty.
        \]
        (Each $\E|X_k| = \E|X_1|$ because the $X_i$ are identically distributed.) $\checkmark$

  \item \textit{Martingale property}: $\E[S_n \mid \F_{n-1}] = S_{n-1}$ a.s.,
        shown above. $\checkmark$
\end{enumerate}

Therefore $(S_n, \F_n)_{n \ge 0}$ is a martingale. \hfill$\square$

\begin{remark}
  The structure here is the template for all random walk martingale arguments:
  \begin{itemize}
    \item The ``known past'' part ($S_{n-1}$) comes out of the expectation via
          measurability (Exercise~3).
    \item The ``new increment'' part ($X_n$) contributes nothing because it is
          independent of the past and has mean zero.
  \end{itemize}
  Keep this two-part decomposition in mind --- it reappears throughout the
  course whenever you need to verify the martingale property for a random walk.
\end{remark}

\end{document}
